\documentclass[aspectratio=169]{beamer}
\usetheme{Madrid}
\usecolortheme{default}

\title[Statistical Depts: GOI \& GOK]{Functioning of Statistical Departments}
\subtitle{Government of India \& Government of Kerala}
\author{Alen Francis Joseph, Semester 1, M.Tech CSE (DS and AI)}
\date{\today}

\begin{document}

\begin{frame}
    \titlepage
\end{frame}

\begin{frame}{Overview of India's Statistical System}
    \begin{itemize}
        \item \textbf{Federal Structure:} India's statistical system operates on a decentralized federal structure, with distinct roles for the central and state governments.
        \item \textbf{Central Nodal Agency:} The Ministry of Statistics and Programme Implementation (MoSPI) acts as the apex body at the national level.
        \item \textbf{State Nodal Agency:} Each state has a State Statistical Bureau (SSB). In Kerala, this is the Department of Economics and Statistics (DES).
        \item \textbf{Objective:} To ensure the availability of reliable, timely, and credible statistics for evidence-based policymaking, economic planning, and governance.
    \end{itemize}
\end{frame}

\begin{frame}{Government of India: MoSPI}
    \begin{itemize}
        \item \textbf{Establishment:} Created in 1999 to oversee official statistics and monitor government programs.
        \item \textbf{Structure:} Divided into two main functional wings:
        \begin{itemize}
            \item \textit{Statistics Wing:} Comprising the National Statistical Office (NSO).
            \item \textit{Programme Implementation Wing:} Monitors central sector infrastructure projects.
        \end{itemize}
        \item \textbf{Mandate:} Sets statistical standards, coordinates statistical work across ministries, and maintains liaisons with international bodies like the UNSD, IMF, and World Bank.
    \end{itemize}
\end{frame}

\begin{frame}{National Statistical Office (NSO) - GoI}
    \begin{itemize}
        \item The NSO is the central repository of statistical data in India.
        \item \textbf{Core Divisions:}
        \begin{itemize}
            \item \textit{Central Statistics Office (CSO):} Responsible for macroeconomic aggregates and national accounting.
            \item \textit{National Sample Survey Office (NSSO):} Conducts large-scale socio-economic surveys across the country.
        \end{itemize}
        \item \textbf{Data Governance:} Works towards improving data quality, transparency, and integrating new digital data sources for better economic assessment.
    \end{itemize}
\end{frame}

\begin{frame}{Key Functions \& Outputs (MoSPI/NSO)}
    \begin{itemize}
        \item \textbf{National Accounts:} Preparation of Gross Domestic Product (GDP), Gross Value Added (GVA), and capital formation estimates.
        \item \textbf{Industrial Statistics:} Compilation and release of the monthly Index of Industrial Production (IIP) and the Annual Survey of Industries (ASI).
        \item \textbf{Price Statistics:} Releases the Consumer Price Index (CPI), which is crucial for measuring inflation and guiding national monetary policy.
        \item \textbf{Socio-Economic Surveys:} Executes the Periodic Labour Force Survey (PLFS) for employment data and the Household Consumption Expenditure Survey (HCES) for poverty estimation.
    \end{itemize}
\end{frame}

\begin{frame}{Government of Kerala: DES Overview}
    \begin{itemize}
        \item \textbf{Nodal Agency:} The Department of Economics and Statistics (DES) functions under the Planning and Economic Affairs Department.
        \item \textbf{Role:} Serves as the primary agency for the collection, compilation, and analysis of statistics relating to all sectors of the Kerala economy.
        \item \textbf{Structure:} Operates through a central directorate in Thiruvananthapuram with extensive district-level and taluk-level survey units.
        \item \textbf{Alignment:} Works in strict coordination with the State Planning Board and local self-government institutions.
    \end{itemize}
\end{frame}

\begin{frame}{Key Functions of Kerala DES}
    \begin{itemize}
        \item \textbf{State Income Estimation:} Calculates Gross State Domestic Product (GSDP) and per capita income specifically for Kerala.
        \item \textbf{Agricultural Statistics:} Conducts the Establishment of an Agency for Reporting Agricultural Statistics (EARAS) to estimate crop area, production, and yield.
        \item \textbf{Vital Statistics:} Compilation of critical data on births, deaths, and demographic trends essential for Kerala's health planning.
        \item \textbf{Local Level Planning:} Provides localized data sets and decentralized planning statistics directly to Gram Panchayats and Municipalities.
    \end{itemize}
\end{frame}

\begin{frame}{Major Publications \& Outputs (Kerala DES)}
    \begin{itemize}
        \item \textbf{Price Statistics:} Collects and publishes daily, weekly, and monthly prices of essential commodities, which drives the state's cost of living index.
        \item \textbf{Industrial Data:} Assists in state-level compilations for the Annual Survey of Industries (ASI) and computes the state-specific IIP.
        \item \textbf{Economic Review:} Contributes extensive primary data for the Kerala State Planning Board's annual "Economic Review" document.
        \item \textbf{Sample Surveys:} Participates actively in NSSO surveys with matching state samples to ensure highly accurate, district-level representation.
    \end{itemize}
\end{frame}

\begin{frame}{Center-State Coordination}
    \begin{itemize}
        \item \textbf{Methodological Synchronization:} MoSPI provides the overarching methodological framework, while the Kerala DES executes state-specific data gathering.
        \item \textbf{Matching Samples:} In NSSO surveys, Kerala DES canvasses a matching "state sample" alongside the GoI's "central sample" to generate robust sub-regional estimates.
        \item \textbf{Capacity Building:} MoSPI offers training and standardizes concepts and definitions to ensure data interoperability across the nation.
        \item \textbf{SDG Monitoring:} Both agencies collaborate closely to monitor India's progress on the UN Sustainable Development Goals (SDGs) at both national and sub-national levels.
    \end{itemize}
\end{frame}

\begin{frame}{Conclusion}
    \begin{itemize}
        \item \textbf{Synergy:} The functioning of India's statistical framework relies heavily on the synergy between national bodies like MoSPI and state bodies like the Kerala DES.
        \item \textbf{Evidence-Based Policy:} From setting national monetary policy to guiding Kerala's decentralized panchayat spending, reliable data forms the backbone of governance.
        \item \textbf{Future Directions:} Both departments are rapidly moving towards digital data collection, real-time analytics, and stronger administrative data integration to reduce survey delays.
    \end{itemize}
\end{frame}

\end{document}
